\section{模块设计}

\subsection{模块划分原则}
本系统的模块划分不按代码目录机械展开，而是按业务职责、数据边界和协作关系进行组织。其核心目标是将当前代码中已经存在的用户、组织、资源、检测、审核、通知和日志能力沉淀为稳定模块，同时为论文检测、评审文本检测和统一报告等迭代需求预留扩展空间。

在划分过程中遵循以下原则：
\begin{itemize}[leftmargin=2em]
    \item \textbf{围绕业务能力划分模块}。模块对应稳定业务职责，而不是简单对应前端页面或后端文件夹；
    \item \textbf{高内聚、低耦合}。单个模块内部尽量完成同一类业务闭环，不直接依赖过多外部模块细节；
    \item \textbf{统一入口、差异化实现}。资源管理、检测任务管理和结果输出尽量统一，不同检测对象的差异通过专用检测模块或适配器处理；
    \item \textbf{兼顾现有实现与后续演进}。模块设计既要能映射当前 Django 后端、Vue 前端、Celery 任务和 AI 服务实现，也要能支撑下一阶段新增论文与评审能力。
\end{itemize}

\subsection{模块总体关系类图}
\begin{figure}[H]
    \centering
    \resizebox{0.98\textwidth}{!}{\begin{tikzpicture}[
    node distance=1.8cm and 1.6cm,
    every node/.style={font=\small},
    umlclass/.style={
        draw,
        rectangle split,
        rectangle split parts=2,
        rectangle split part align={center,left},
        rounded corners=3pt,
        line width=0.8pt,
        minimum width=4.2cm,
        text width=4.2cm,
        inner sep=5pt
    },
    relation/.style={-{Latex[length=2.3mm]}, thick, draw=black!75}
]

\node[umlclass, fill=blue!8] (auth) {
\textbf{用户与权限模块}
\nodepart{second}
+ 注册登录\\
+ 角色权限\\
+ 身份认证
};

\node[umlclass, fill=blue!8, right=of auth] (org) {
\textbf{组织与成员模块}
\nodepart{second}
+ 组织申请\\
+ 成员关系\\
+ 组织配额
};

\node[umlclass, fill=green!10, below=of auth] (resource) {
\textbf{资源管理模块}
\nodepart{second}
+ 文件上传\\
+ 资源登记\\
+ PDF图片提取
};

\node[umlclass, fill=orange!12, right=of resource] (task) {
\textbf{检测任务管理模块}
\nodepart{second}
+ 任务创建\\
+ 状态流转\\
+ 结果聚合
};

\node[umlclass, fill=orange!8, below left=of task] (image) {
\textbf{图像检测模块}
\nodepart{second}
+ 图像伪造识别\\
+ 子方法结果\\
+ 可疑区域生成
};

\node[umlclass, fill=orange!8, below=of task] (paper) {
\textbf{论文检测模块}
\nodepart{second}
+ 论文解析\\
+ 段落切分\\
+ AIGC分析
};

\node[umlclass, fill=orange!8, below right=of task] (review) {
\textbf{评审检测模块}
\nodepart{second}
+ 模板化分析\\
+ 生成痕迹识别\\
+ 评审结果封装
};

\node[umlclass, fill=purple!8, right=of task] (manual) {
\textbf{人工审核模块}
\nodepart{second}
+ 审核申请\\
+ 审核分派\\
+ 审核结论汇总
};

\node[umlclass, fill=green!10, below=of manual] (report) {
\textbf{报告服务模块}
\nodepart{second}
+ 检测报告\\
+ 人工审核报告\\
+ 文件归档
};

\node[umlclass, fill=yellow!15, left=of report] (notify) {
\textbf{通知消息模块}
\nodepart{second}
+ 站内通知\\
+ WebSocket推送\\
+ 状态反馈
};

\node[umlclass, fill=gray!15, left=of image] (log) {
\textbf{统计日志模块}
\nodepart{second}
+ 操作日志\\
+ 任务统计\\
+ 审核统计
};

\draw[relation] (auth) -- node[above, font=\scriptsize, fill=white, inner sep=1pt] {组织归属} (org);
\draw[relation] (auth) -- node[left, font=\scriptsize, fill=white, inner sep=1pt] {权限控制} (resource);
\draw[relation] (org) -- node[right, font=\scriptsize, fill=white, inner sep=1pt] {成员约束} (task);
\draw[relation] (resource) -- node[above, font=\scriptsize, fill=white, inner sep=1pt] {创建任务} (task);
\draw[relation] (task) -- node[above left, font=\scriptsize, fill=white, inner sep=1pt] {调度} (image);
\draw[relation] (task) -- node[right, font=\scriptsize, fill=white, inner sep=1pt] {调度} (paper);
\draw[relation] (task) -- node[above right, font=\scriptsize, fill=white, inner sep=1pt] {调度} (review);
\draw[relation] (task) -- node[above, font=\scriptsize, fill=white, inner sep=1pt] {申请复核} (manual);
\draw[relation] (manual) -- node[right, font=\scriptsize, fill=white, inner sep=1pt] {生成} (report);
\draw[relation] (task) -- node[below, font=\scriptsize, fill=white, inner sep=1pt] {生成} (report);
\draw[relation] (task) -- node[above, font=\scriptsize, fill=white, inner sep=1pt] {通知} (notify);
\draw[relation] (manual) -- node[below, font=\scriptsize, fill=white, inner sep=1pt] {通知} (notify);
\draw[relation] (log) -- node[below, font=\scriptsize, fill=white, inner sep=1pt] {记录} (resource);
\draw[relation] (log) -- node[left, font=\scriptsize, fill=white, inner sep=1pt] {记录} (task);
\draw[relation] (log) -- node[above, font=\scriptsize, fill=white, inner sep=1pt] {记录} (manual);

\end{tikzpicture}
}
    \caption{系统模块关系 UML 类图}
\end{figure}

图 4-1 展示了本系统在模块层级上的主要组成和依赖关系。用户与权限模块、组织与成员模块和资源管理模块构成平台基础能力；检测任务管理模块作为任务中枢，统一调度图像检测、论文检测和评审检测三个专业检测模块；自动检测结果可进一步进入人工审核模块进行复核；报告服务模块对检测和审核结果进行汇总输出；通知消息模块和统计日志模块则贯穿全流程，承担状态反馈与治理支撑职责。

\subsection{模块总览}
\begin{longtable}{|p{3.2cm}|p{6.7cm}|p{4.1cm}|}
\hline
模块名称 & 主要职责 & 主要依赖 \\ \hline
用户与权限模块 & 负责注册登录、身份认证、角色区分、权限校验和账户状态管理 & 组织与成员模块、通知消息模块 \\ \hline
组织与成员模块 & 负责组织申请、组织信息维护、邀请码管理、成员关系和组织级配额管理 & 用户与权限模块、统计日志模块 \\ \hline
资源管理模块 & 负责文件上传、资源登记、分类标记、图片提取和资源查询 & 用户与权限模块、检测任务管理模块、文件存储 \\ \hline
检测任务管理模块 & 负责任务创建、资源绑定、状态流转、异步调度、结果聚合和任务级报告关联 & 资源管理模块、各检测模块、报告服务模块、通知消息模块 \\ \hline
图像检测模块 & 负责图像预处理、图像真伪检测、子方法结果汇总和可疑区域生成 & 检测任务管理模块、AI 推理服务 \\ \hline
论文检测模块 & 负责论文解析、文本切分、AIGC 风险分析和论文结果封装 & 检测任务管理模块、AI 推理服务 \\ \hline
评审检测模块 & 负责评审文本分析、模板化倾向识别和生成式痕迹分析 & 检测任务管理模块、AI 推理服务 \\ \hline
人工审核模块 & 负责审核申请、审核分派、人工评分、审核结论汇总和人工审核报告生成 & 检测任务管理模块、用户与权限模块、通知消息模块 \\ \hline
报告服务模块 & 负责任务报告和人工审核报告生成、下载与归档 & 检测任务管理模块、人工审核模块、文件存储 \\ \hline
通知消息模块 & 负责站内通知、实时推送、状态更新提醒和审核结果通知 & 用户与权限模块、检测任务管理模块、人工审核模块 \\ \hline
统计日志模块 & 负责操作日志、任务统计、审核统计和后台治理支撑 & 全部业务模块 \\ \hline
\end{longtable}

\subsection{用户与权限模块}
用户与权限模块是整个平台的访问控制入口，负责用户注册、登录、身份认证、角色管理和基础权限校验。结合当前后端实现，该模块主要围绕 \texttt{User} 模型展开，并通过角色字段与权限位控制发布者、审核者和管理员在上传、提交、发布和审核等操作上的差异。

该模块的主要输入包括注册信息、登录凭证、邀请码和用户会话请求，主要输出包括认证结果、用户身份信息、角色信息和权限判断结果。由于用户权限与组织归属密切相关，因此该模块与组织与成员模块存在紧密协作关系。

从迭代角度看，该模块后续还需要继续承担统一权限入口职责，即论文检测、评审检测、模型配置和综合治理等新功能都不应绕过该模块单独实现权限控制。

\subsection{组织与成员模块}
组织与成员模块负责组织生命周期和组织内部协作关系管理，包括组织申请、组织审核、组织信息维护、管理员归属、邀请码分发、发布者与审核者关系维护，以及组织级配额管理等功能。当前代码中的 \texttt{OrganizationApplication}、\texttt{Organization}、\texttt{InvitationCode} 和 \texttt{PublisherReviewerRelationship} 等实体已经为该模块提供了实现基础。

该模块的输入主要包括组织申请资料、组织基本信息、成员关系调整请求和邀请码使用请求，输出则包括组织状态、成员归属关系和组织级资源使用能力。由于系统的用户角色和审核协作链条都依赖组织结构，因此该模块是平台治理能力的基础。

后续若要支持更复杂的机构治理，例如多级管理员、模型使用额度配置和组织维度统计看板，也应继续在该模块上扩展，而不是散落到各业务模块中实现。

\subsection{资源管理模块}
资源管理模块负责系统中所有待检测内容的接入、登记和管理，是检测业务链路的起点。当前实现已经围绕文件管理、图片上传和 PDF 中图片提取建立了较完整的图像资源处理链路，后续则需要把这一能力统一提升为面向图像、论文和评审文本三类资源的资源接入能力。

在当前代码中，该模块可映射到 \texttt{FileManagement}、\texttt{ImageUpload} 及相关上传接口和资源查询接口。其核心职责包括：
\begin{itemize}[leftmargin=2em]
    \item 接收用户上传的原始文件或文本；
    \item 记录资源元数据、分类标签、来源信息和所属组织；
    \item 对 PDF 类资源执行图片提取等预处理；
    \item 为后续检测任务提供统一的资源引用和状态基础。
\end{itemize}

资源管理模块对外输出的是“可被检测任务引用的资源对象”，而不是直接输出检测结果。这一点有助于保持上传流程与检测流程的分离，使后续支持资源复检、批量检测和多类型资源扩展更容易实现。

\subsection{检测任务管理模块}
检测任务管理模块是整个系统的业务中枢，负责将资源接入、检测执行、状态流转、结果聚合和报告输出串联为一条完整链路。当前代码中的 \texttt{DetectionTask}、\texttt{DetectionResult}、\texttt{SubDetectionResult}，以及 Celery 任务和任务状态通知逻辑，共同构成了该模块的实现基础。

该模块的主要职责包括：
\begin{itemize}[leftmargin=2em]
    \item 创建检测任务并绑定待检测资源；
    \item 初始化任务状态、检测方式和报告占位信息；
    \item 根据任务类型将请求投递到异步任务系统；
    \item 汇总单资源结果和子方法结果；
    \item 在任务完成后触发报告生成与通知发送。
\end{itemize}

在后续迭代中，该模块需要从“图像任务中心”提升为“统一检测任务中心”。也就是说，无论是图像检测、论文检测还是评审检测，都应尽量通过该模块进入统一状态机和统一任务生命周期，而不在外围额外构建平行任务体系。

\subsection{图像检测模块}
图像检测模块是当前系统中实现最成熟的专业检测模块，负责学术图像伪造识别、多检测方法执行、可疑区域生成和图像级结果输出。当前后端通过任务模块、AI 服务调用逻辑和 \texttt{SubDetectionResult} 子结果结构，已经具备图像多方法综合检测基础。

该模块的主要输入是图像资源、图像任务上下文和检测配置，输出则包括图像级真假判断、置信度、子方法概率、掩码图、可视化产物和图像级说明信息。由于图像检测往往耗时较长，因此它不直接面向前端同步返回，而是通过检测任务管理模块与异步任务系统协同完成执行。

从设计上看，该模块是后续其他检测模块的参考模板。论文检测和评审检测虽然算法不同，但在“接收任务上下文、执行检测、回写结构化结果”这一模式上可以沿用相同接口思想。

\subsection{论文检测模块}
论文检测模块是本轮迭代新增的重点模块之一，负责对整篇论文进行解析、切分和 AIGC 风险分析，并输出论文级别和段落级别的检测结果。当前代码中尚未形成完整实现，因此本章对该模块主要给出目标结构设计。

该模块的主要职责包括：
\begin{itemize}[leftmargin=2em]
    \item 解析论文文件内容并提取结构化文本；
    \item 按章节、段落或语义单元进行文本切分；
    \item 调用对应模型评估生成式痕迹或文本异常模式；
    \item 将段落级检测结果汇总为论文级风险结论；
    \item 为报告服务模块提供可直接引用的结构化结果。
\end{itemize}

该模块不应复用图像检测模块中的结果结构细节，而应在统一任务框架下采用专属结果载荷。这样既能共享任务生命周期，又能避免论文结果被迫适配为图像式结构。

\subsection{评审检测模块}
评审检测模块面向 Review 文本，重点关注模板化倾向、重复性表达、生成式语言痕迹和不自然的学术评审模式。该模块与论文检测模块同属文本类检测能力，但其输入对象、特征关注点和输出形式不同，因此应作为独立模块设计，而不是简单并入论文检测模块。

该模块的主要输入是评审文本和任务上下文，输出包括模板化风险、生成式风险、可疑片段说明和评审级结论。其在系统中的位置与图像检测模块、论文检测模块并列，统一由检测任务管理模块调度。

这样设计的好处是：未来如果评审检测引入更细粒度规则引擎、相似评审比对或多轮评审一致性分析，不需要对论文检测模块进行破坏性改造。

\subsection{人工审核模块}
人工审核模块负责对自动检测结果进行人工复核，是保证平台可信度和可解释性的关键模块。当前代码已经围绕 \texttt{ReviewRequest}、\texttt{ManualReview}、\texttt{ImageReview} 和审核相关视图接口实现了较完整的图像人工审核链路。

该模块的主要职责包括：
\begin{itemize}[leftmargin=2em]
    \item 接收用户发起的人工审核申请；
    \item 由管理员审批并分派审核人员；
    \item 支持审核人员逐项查看自动检测结果并提交人工判断；
    \item 汇总多名审核人员意见形成审核结论；
    \item 在审核完成后生成人工审核报告并反馈给申请方。
\end{itemize}

从演进角度看，该模块后续不应只服务图像资源，还应逐步扩展为“统一人工复核模块”。也就是说，论文和评审文本在必要时也应能纳入人工审核流程，只是其评分维度和审核界面可能与图像不同。

\subsection{报告服务模块}
报告服务模块负责将检测结果和人工审核结果整理为对用户可见、可下载、可归档的正式输出。当前后端已经通过 \texttt{report\_generator.py} 具备任务检测报告和人工审核报告的 PDF 生成能力，因此该模块已有较明确的实现基础。

该模块的主要输入包括任务级结果、子检测结果、人工审核结论和基础元数据，输出包括任务报告、单资源报告、人工审核报告以及后续可能扩展的综合报告。由于报告本质上是对其他模块输出的汇总，因此该模块不负责产生原始检测结论，而负责对结论进行结构化表达、格式化展示和文件归档。

在后续迭代中，该模块需要从“图像任务报告生成器”演进为“统一报告服务”，使图像、论文、评审和人工审核结果都能通过统一模板策略输出。

\subsection{通知消息模块}
通知消息模块负责向用户和管理员反馈系统中的关键状态变化，包括任务开始、任务完成、报告生成、审核申请、审核结论和系统公告等。当前实现已经通过 \texttt{Notification} 模型、通知接口和 Django Channels WebSocket 消费者支持实时消息推送。

该模块的价值不在于保存单条消息本身，而在于作为系统状态外显层，让异步任务和长时审核流程能够及时反馈给前端用户。对于以异步检测为核心的平台而言，通知模块是体验闭环的一部分，而不是可有可无的附属功能。

\subsection{统计日志模块}
统计日志模块负责记录平台运行过程中的关键操作和业务指标，为管理端统计分析、问题追踪和后续审计提供支撑。当前代码中的 \texttt{Log} 模型、管理端统计接口和部分看板接口已构成该模块的基础。

该模块的主要职责包括：
\begin{itemize}[leftmargin=2em]
    \item 记录上传、检测、审核申请、人工审核等关键操作日志；
    \item 支持按用户、组织、任务和时间范围进行查询；
    \item 为管理端提供任务统计、审核统计和趋势分析数据；
    \item 为后续论文日志、评审日志和模型使用日志扩展提供统一治理入口。
\end{itemize}

考虑到本项目后续治理能力会不断增强，日志与统计模块不应被视为简单的后台附属页面，而应作为平台可追溯性和可管理性的基础模块持续完善。

\subsection{模块协作关系说明}
从整体协作看，系统模块之间形成了较清晰的主干链路。用户与权限模块和组织与成员模块提供访问与归属基础；资源管理模块负责将检测对象接入平台；检测任务管理模块负责统一编排自动检测流程；图像检测、论文检测和评审检测模块分别提供专业检测能力；人工审核模块在需要时对自动结果进行补充复核；报告服务模块负责输出正式文档；通知消息模块和统计日志模块则在全过程中提供反馈和治理支撑。

这一协作结构的优势在于：
\begin{itemize}[leftmargin=2em]
    \item 主业务链条清晰，便于后续逐步实现论文检测和评审检测；
    \item 模块间职责边界相对稳定，降低了“功能增长即结构失控”的风险；
    \item 与当前代码已有实现具有较强映射关系，便于后续详细设计与实际开发衔接。
\end{itemize}

\subsection{本章小结}
本章从业务职责角度对系统进行了模块化拆分，并结合当前代码结构与后续迭代目标，给出了模块边界、依赖关系和演进方向。后续详细设计可在本章基础上进一步细化每个模块的实体结构、接口定义、时序流程和异常处理策略。
